% Options for packages loaded elsewhere
\PassOptionsToPackage{unicode}{hyperref}
\PassOptionsToPackage{hyphens}{url}
\PassOptionsToPackage{dvipsnames,svgnames,x11names}{xcolor}
\documentclass[
  10pt,
  fontset=fandol]{ctexart}
\usepackage{xcolor}
\usepackage[margin=16mm]{geometry}
\usepackage{amsmath,amssymb}
\setcounter{secnumdepth}{-\maxdimen} % remove section numbering
\usepackage{iftex}
\ifPDFTeX
  \usepackage[T1]{fontenc}
  \usepackage[utf8]{inputenc}
  \usepackage{textcomp} % provide euro and other symbols
\else % if luatex or xetex
  \usepackage{unicode-math} % this also loads fontspec
  \defaultfontfeatures{Scale=MatchLowercase}
  \defaultfontfeatures[\rmfamily]{Ligatures=TeX,Scale=1}
\fi
\usepackage{lmodern}
\ifPDFTeX\else
  % xetex/luatex font selection
\fi
% Use upquote if available, for straight quotes in verbatim environments
\IfFileExists{upquote.sty}{\usepackage{upquote}}{}
\IfFileExists{microtype.sty}{% use microtype if available
  \usepackage[]{microtype}
  \UseMicrotypeSet[protrusion]{basicmath} % disable protrusion for tt fonts
}{}
\makeatletter
\@ifundefined{KOMAClassName}{% if non-KOMA class
  \IfFileExists{parskip.sty}{%
    \usepackage{parskip}
  }{% else
    \setlength{\parindent}{0pt}
    \setlength{\parskip}{6pt plus 2pt minus 1pt}}
}{% if KOMA class
  \KOMAoptions{parskip=half}}
\makeatother
\setlength{\emergencystretch}{3em} % prevent overfull lines
\providecommand{\tightlist}{%
  \setlength{\itemsep}{0pt}\setlength{\parskip}{0pt}}
\usepackage{amsmath,amssymb,mathtools,needspace}
\xeCJKDeclareCharClass{CJK}{"2032,"0400->"04FF,"2160->"216B,"2460->"2473,"25A0,"25C7,"25CB,"25CF}
\IfFontExistsTF{Arial Unicode MS}{\xeCJKsetup{AutoFallBack=true}\setCJKfallbackfamilyfont{\CJKrmdefault}{Arial Unicode MS}}{}
\setcounter{secnumdepth}{0}
\setlength{\emergencystretch}{2em}
\allowdisplaybreaks
\usepackage{bookmark}
\IfFileExists{xurl.sty}{\usepackage{xurl}}{} % add URL line breaks if available
\urlstyle{same}
\hypersetup{
  pdftitle={一阶方程组},
  colorlinks=true,
  linkcolor={Maroon},
  filecolor={Maroon},
  citecolor={Blue},
  urlcolor={Blue},
  pdfcreator={LaTeX via pandoc}}

\title{一阶方程组}
\author{}
\date{}

\begin{document}
\maketitle

\subsubsection{5.6.1 一阶方程组}\label{ux4e00ux9636ux65b9ux7a0bux7ec4}

前面研究了单个方程 \(y'=f\) 的数值解法，只要把 \(y\) 和 \(f\) 理解为向量，那么，所提供的各种计算格式即可应用到一阶方程组的情形。

考察一阶方程组

\[
y'_i=f_i(x,y_1,y_2,\cdots,y_N)\qquad(i=1,2,\cdots,N)
\]

的初值问题，初始条件为

\[
y_i(x_0)=y_i^0\qquad(i=1,2,\cdots,N).
\]

若采用向量的记号，记

\[
\boldsymbol y=(y_1,y_2,\cdots,y_N)^{\mathrm T},\qquad
\boldsymbol y^0=(y_1^0,y_2^0,\cdots,y_N^0)^{\mathrm T},
\]

\[
\boldsymbol f=(f_1,f_2,\cdots,f_N)^{\mathrm T},
\]

则上述方程组的初值问题可表示为 \(\begin{cases}\boldsymbol y'=\boldsymbol f(x,\boldsymbol y),\\\boldsymbol y(x_0)=\boldsymbol y_0.\end{cases}\) 求解这一初值问题的四阶 Runge-Kutta 格式为

\[
\boldsymbol y_{n+1}=\boldsymbol y_n+\frac h6(\boldsymbol k_1+2\boldsymbol k_2+2\boldsymbol k_3+\boldsymbol k_4).
\]

其中

\[
\boldsymbol k_1=\boldsymbol f(x_n,\boldsymbol y_n),\qquad
\boldsymbol k_2=\boldsymbol f\left(x_n+\frac h2,\boldsymbol y_n+\frac h2\boldsymbol k_1\right),
\]

\[
\boldsymbol k_3=\boldsymbol f\left(x_n+\frac h2,\boldsymbol y_n+\frac h2\boldsymbol k_2\right),\qquad
\boldsymbol k_4=\boldsymbol f(x_n+h,\boldsymbol y_n+h\boldsymbol k_3);
\]

或表示为

\[
y_{i,n+1}=y_{in}+\frac h6(K_{i1}+2K_{i2}+2K_{i3}+K_{i4})\qquad(i=1,2,\cdots,N),
\]

其中

\[
\begin{aligned}
K_{i1}&=f_i(x_n,y_{1n},y_{2n},\cdots,y_{Nn}),\\
K_{i2}&=f_i\left(x_n+\frac h2,y_{1n}+\frac h2K_{11},y_{2n}+\frac h2K_{21},\cdots,y_{Nn}+\frac h2K_{N1}\right),\\
K_{i3}&=f_i\left(x_n+\frac h2,y_{1n}+\frac h2K_{12},y_{2n}+\frac h2K_{22},\cdots,y_{Nn}+\frac h2K_{N2}\right),\\
K_{i4}&=f_i(x_n+h,y_{1n}+hK_{13},y_{2n}+hK_{23},\cdots,y_{Nn}+hK_{N3}).
\end{aligned}
\]

这里 \(y_{in}\) 是第 \(i\) 个因变量 \(y_i(x)\) 在节点 \(x_n=x_0+nh\) 的近似值。

为了帮助理解这一格式的计算过程，再考察以下两个方程的特殊情形：

\href{../../source-images/pdf-146.jpeg}{原页}

\[
\begin{cases}
y'=f(x,y,z),\quad y(x_0)=y_0,\\
z'=g(x,y,z),\quad z(x_0)=z_0,
\end{cases}
\]

这时四阶 Runge-Kutta 格式具有如下形式：

\[
\begin{cases}
\displaystyle y_{n+1}=y_n+\frac h6(K_1+2K_2+2K_3+K_4),\\[6pt]
\displaystyle z_{n+1}=z_n+\frac h6(L_1+2L_2+2L_3+L_4),
\end{cases}
\tag{5.6.1}
\]

其中

\[
\begin{cases}
K_1=f(x_n,y_n,z_n),\\
\displaystyle K_2=f\left(x_n+\frac h2,y_n+\frac h2K_1,z_n+\frac h2L_1\right),\\[4pt]
\displaystyle K_3=f\left(x_n+\frac h2,y_n+\frac h2K_2,z_n+\frac h2L_2\right),\\[4pt]
K_4=f(x_n+h,y_n+hK_3,z_n+hL_3),\\
L_1=g(x_n,y_n,z_n),\\
\displaystyle L_2=g\left(x_n+\frac h2,y_n+\frac h2K_1,z_n+\frac h2L_1\right),\\[4pt]
\displaystyle L_3=g\left(x_n+\frac h2,y_n+\frac h2K_2,z_n+\frac h2L_2\right),\\[4pt]
L_4=g(x_n+h,y_n+hK_3,z_n+hL_3).
\end{cases}
\tag{5.6.2}
\]

这是一步法，利用节点 \(x_n\) 上的值 \(y_n,z_n\)，由式（5.6.2）顺序计算 \(K_1,\allowbreak L_1,\allowbreak K_2,\allowbreak L_2,\allowbreak K_3,\allowbreak L_3,\allowbreak K_4,\allowbreak L_4\)，然后代入式（5.6.1）即可求得节点 \(x_{n+1}\) 上的 \(y_{n+1},z_{n+1}\)。

\end{document}
